\paragraph{C.3k (Proved): torsion injectivity in dimension three and the ordinary algebraic-equivalence consequence.}
The cyclic-cover mechanism used in
\citet[Section 2.5, p.~466]{D_Kollar2016LocalPicard}
also gives the following precise statement in dimension
three. We give the normalization, unit and connectedness
arguments, since they are indispensable at this boundary.

\emph{(i) Prime-to-residue-characteristic torsion.}
Let $(R,\mathfrak m,k)$ be an excellent, normal, strictly
henselian local ring of dimension $N\ge3$. Let
$0\ne f\in\mathfrak m$, and put
\[
 X=\operatorname{Spec}R\setminus\{\mathfrak m\},
 \qquad Y=X\cap V(f).
\]
Assume that $Y$ is regular. For every positive integer
$n$ invertible in $R$, restriction induces an injection
\begin{equation}\label{eq:C3-prime-to-p-Pic-injection}
              \operatorname{Pic}(X)[n]
                    \longrightarrow\operatorname{Pic}(Y)[n].
\end{equation}
Here the line bundles and the assertion are integral.
The result applies in equal and mixed characteristic,
and does not require algebraizability. Normality of
$R$ suffices on the source; regularity of $X$ is an
allowed stronger hypothesis from the manuscript.

\begin{proof}
We first determine the global units that enter Kummer
theory. Since $R$ is normal and $N\ge2$, the local
Hartogs property gives
\[
                  \Gamma(X,\mathcal O_X)=R.
\]
For completeness, the extension statement follows by
expressing a normal noetherian domain as the intersection
of its height-one localizations inside its fraction
field: every height-one point lies in $X$.
Every unit of $R$ has an $n$th root. Its reduction has
a root in the separably closed residue field, and
the derivative of $T^n-u$ is invertible at that root;
the henselian property supplies the lift.

We need the corresponding statement on $Y$, where
$R/(f)$ need not be normal or $S_2$. Put $C=R/(f)$.
This ring is reduced. Indeed normality gives
$\operatorname{depth}R\ge2$, and $f$ is a non-zero-divisor,
so $\operatorname{depth}C\ge1$. At every point except
its closed point, $C$ is regular by hypothesis.
It therefore satisfies $S_1$ everywhere and is
generically reduced, which together imply reducedness.
Each component has dimension $N-1\ge2$: the minimal
primes over $(f)$ have height one, and excellence
gives the catenary dimension formula in the local
domain $R$.

Let $\widetilde C$ be the normalization of $C$ in
its total ring of fractions. It is finite over $C$
because $C$ is excellent, hence Nagata; see
\citet[Lemma 29.54.14, Tag 0AVK]{D_StacksLocalPicard}.
Normalization is an isomorphism above the regular
open $Y$. A finite algebra over a henselian local
ring is a product of henselian local rings
\citep[Lemma 10.153.4, Tag 04GH]{D_StacksLocalPicard},
so
\[
                    \widetilde C=\prod_{a=1}^{s} C_a,
\]
where each $C_a$ is a normal henselian local domain
of dimension $N-1$. Its residue field is a finite
extension of the separably closed field $k$, hence
is purely inseparable and still separably closed.
Thus each $C_a$ is strictly henselian. Removing the
closed points of $\operatorname{Spec}\widetilde C$
gives exactly $Y$. Applying the same height-one
intersection argument to the normal rings $C_a$
gives the equality, not merely an inclusion,
\begin{equation}\label{eq:C3-section-finite-normalization}
             \Gamma(Y,\mathcal O_Y)=\widetilde C.
\end{equation}
In particular this ring is finite over $C$, and
its group of units is $n$-divisible, by Hensel's
lemma in each factor. This supplies a global
$n$th root for every unit that can occur in a
trivialization on $Y$.

The Kummer sequences on $X$ and $Y$, together with
these two divisibility statements, now give
compatible isomorphisms
\begin{equation}\label{eq:C3-Kummer-Pic-identifications}
\begin{aligned}
 H^1_{\mathrm{et}}(X,\mu_n)&\simeq\operatorname{Pic}(X)[n],\\
 H^1_{\mathrm{et}}(Y,\mu_n)&\simeq\operatorname{Pic}(Y)[n].
\end{aligned}
\end{equation}
Suppose there were a nontrivial kernel element of
order $d$ dividing $n$. Use $d$ in place of $n$.
All $d$th roots of unity lie in $R$, so $\mu_d$ is
the constant cyclic group on the connected normal
scheme $X$. Its Kummer class has exact order $d$.
Under
\[
 H^1_{\mathrm{et}}(X,\mu_d)
       =\operatorname{Hom}_{\mathrm{cont}}
                       (\pi_1^{\mathrm{et}}(X),\mu_d),
\]
the order of a character equals the order of its
image. Thus this class defines a connected finite
etale $\mu_d$-torsor $X'\to X$ of degree $d$.
Its restriction to $Y$ is the trivial torsor by
\eqref{eq:C3-Kummer-Pic-identifications}. In particular
\begin{equation}\label{eq:C3-split-restricted-cover}
                       X'\times_XY\simeq\coprod_{j=1}^{d}Y.
\end{equation}

The connected normal scheme $X'$ is integral.
Let $S$ be the normalization of $R$ in its function
field. This is a finite normal $R$-algebra, again
by excellence and the finite field extension.
As a finite algebra over the henselian ring $R$,
it is a product of local rings; since it is a
domain, there is only one factor. Consequently
$S$ is a normal henselian local domain, of
dimension $N$. Its punctured spectrum is $X'$:
both are the normalization of $X$ in the same
field. There is one closed point above
$\mathfrak m$, and therefore
\[
 X'\times_XY
       =\operatorname{Spec}(S/fS)\setminus\{\mathfrak m_S\}.
\]

We show that this punctured divisor is connected.
The ring $S$ is excellent. Its completion
$\widehat S$ is a normal local domain
\citep[Lemma 15.53.6, Tag 0C23]{D_StacksLocalPicard}.
The spectrum of this complete domain is irreducible
of dimension $N$, and
$\operatorname{cd}(\widehat S,(f))\le1$, as is seen
from the two-term Cech complex for a principal
ideal. The connectedness theorem
\citep[Theorem 1.6, p.~2497]{D_Varbaro2009},
in the equivalent punctured form of
\citet[Lemma 52.14.1, Tag 0ECR]{D_StacksLocalPicard},
applies because
\[
              1<N-1=\min\{N,N-1\}.
\]
It gives that
$\operatorname{Spec}(\widehat S/f\widehat S)
\setminus\{\mathfrak m_{\widehat S}\}$ is connected.
Completion is faithfully flat and its fibre at
the closed point is the residue field. Hence
the induced morphism from this puncture onto
$\operatorname{Spec}(S/fS)\setminus\{\mathfrak m_S\}$
is surjective. A continuous image of a connected
space is connected, proving the claim before
completion as well. Since $Y$ is nonempty and
$d>1$, this contradicts
\eqref{eq:C3-split-restricted-cover}.
Thus the proposed torsion kernel is zero.
\end{proof}

\emph{(ii) Ordinary algebraic equivalence in
algebraizable characteristic zero.}
Let $k$ be algebraically closed of characteristic
zero, let $A$ be a normal local ring essentially
of finite type over $k$ with residue field $k$,
and put $R=A^{\mathrm{sh}}$. In the setting above,
assume also that $X$ is regular. Then for every
$N\ge3$ the restriction
\begin{equation}\label{eq:C3-ordinary-alg-N3}
          \operatorname{Pic}(X)_{\mathrm{alg}}
                  \longrightarrow\operatorname{Pic}(Y)_{\mathrm{alg}}
\end{equation}
is injective integrally, where algebraic equivalence
is defined by families on the fixed scheme times
a smooth connected parameter curve. It follows
that the rational ordinary $J^1$ restriction is
injective in this exact scope, including $N=3$.

\begin{proof}
After descent of $f$ to one pointed etale
neighbourhood, use that neighbourhood as the
algebraic germ; this leaves $R$ unchanged.
The source local Picard scheme $G$ exists for
the normal algebraic germ of dimension at least
three, with $G(k)=\operatorname{Pic}(X)$.
The finite family-descent, smooth projective
compactification and Poincare-family argument
in C.3j uses this source representability and
characteristic-zero resolution, but does not
use $N\ge4$ or unipotence. Applied here, it shows
that every ordinarily algebraically trivial
line bundle belongs to an abelian subvariety
$B\subset G$.

The kernel of restriction on $B$ is a closed
subgroup scheme of $B$, including when $N=3$.
To justify this without assuming a finite-type
Picard scheme for a two-dimensional target,
take the finite normalization of the algebraic
hyperplane germ. It is an isomorphism away
from the closed point because that puncture
is regular. Its finitely many local factors
are normal of dimension $N-1\ge2$, and hence
have depth at least two. The unit section of
each local Picard functor is represented by
a finitely presented closed immersion by
\citet[Theorem 1(i), Introduction, p.~VII]{D_Boutot1978};
the exact statement is also recorded in
\citet[Definition 3.3, p.~468]{D_Kollar2016LocalPicard}.
Pulling back the product of these unit
sections along the morphism of group functors
\[
 B\lhook\joinrel\longrightarrow G
       \longrightarrow
       \prod_a\mathbf{Pic}^{\mathrm{loc}}(y_a,\widetilde D_a)
\]
gives a closed subgroup scheme $K_B\subset B$.
Here $\widetilde D_a$ are the normal local
factors of the finite normalization and
$y_a$ their closed points; the second arrow
is restriction of local Picard functors.
This construction does not require a
universal line bundle on $B$.
Its $k$-points are precisely
the line bundles in $B(k)$ trivial on $Y$.
Etale base change of the finite normalization
and then the strict henselization preserve
its puncture, so this describes the required
restriction, not a different algebraic model.

If $K_B$ had a nonidentity $k$-point, it would
have a nonidentity torsion $k$-point. Indeed
its reduced identity component is an abelian
subvariety: if it has positive dimension,
its $\ell$-torsion is nontrivial for any prime
$\ell$; if its dimension is zero, the group
is finite and every $k$-point is torsion.
Here group schemes of finite type in
characteristic zero are smooth, so no
infinitesimal qualification is needed.
The inclusion $B\subset G$ would make this
a nontrivial torsion line bundle on $X$
restricting trivially to $Y$, contrary to
part (i), since every positive integer is
invertible in characteristic zero. Thus
$K_B(k)=\{1\}$, proving
\eqref{eq:C3-ordinary-alg-N3}.
\end{proof}

Part (ii) strengthens the injectivity statement
of C.3j, but does not assert surjectivity in
dimension at least four. Part (i) is valid
for arbitrary excellent normal strictly
henselian rings in the stated characteristics;
the ordinary-algebraic-equivalence conclusion
uses the separate algebraizable
characteristic-zero hypotheses explicitly.
